\section{Interpretación de los resultados: mejora del desempeño, brecha entre modelos y costo computacional}

\subsection{De una línea base baja a un nivel cercano al humano}
El expositor mostró que, mediante mejores estrategias de razonamiento y mecanismos de evaluación, el desempeño del sistema puede mejorar de manera notable a partir de un nivel bajo, y en algunas configuraciones acercarse al desempeño humano. Este resultado muestra que el «diseño del sistema» es igual de decisivo para la mejora del desempeño.

\begin{figure}[H]
\centering
\includegraphics[width=0.82\textwidth]{frame-04-human-level.jpg}
\caption{Video aproximadamente en 00:22:23: el expositor muestra la mejora del desempeño y menciona el rango cercano al human level}
\end{figure}

\begin{importantbox}{El verdadero significado de estos resultados}
No significa que el Agente «ya haya resuelto la automatización web», sino que muestra que, dentro del marco correcto de evaluación y búsqueda, las tareas de cadena larga tienen un margen claro de optimización.
\end{importantbox}

\subsection{Evolución entre generaciones de modelos}
El expositor mencionó la tendencia de mejora de Gemini Pro, GPT-4 y modelos posteriores en este tipo de tareas. La capacidad de los modelos en efecto ha avanzado, pero el sistema de apoyo de «planear-ejecutar-verificar» sigue siendo una condición necesaria.

\begin{figure}[H]
\centering
\includegraphics[width=0.82\textwidth]{frame-05-model-comparison.jpg}
\caption{Video aproximadamente en 00:25:15: el expositor discute la mejora en el desempeño de las tareas entre Gemini Pro, GPT-4 y modelos posteriores}
\end{figure}

\begin{knowledgebox}{Cómo interpretar la mejora de los modelos}
\begin{itemize}
    \item Los modelos más fuertes reducen la tasa de error por paso;
    \item pero la estabilidad de la cadena larga sigue dependiendo del mecanismo de búsqueda y retroceso;
    \item mientras más fuerte es el modelo, más se necesita una evaluación y una gobernanza más estrictas, para evitar indicadores infladamente altos.
\end{itemize}
\end{knowledgebox}

\subsection{Equilibrio entre costo y beneficio}
La búsqueda y la verificación elevan la tasa de éxito, pero también aumentan el costo de tokens y de llamadas al entorno. Un sistema en producción debe gestionar de forma explícita el equilibrio entre «el beneficio de la tasa de éxito y el gasto en costo».

\begin{warningbox}{Errores comunes en la gestión de costos}
\begin{itemize}
    \item Perseguir solo la tasa de éxito, sin considerar el costo por tarea;
    \item pasar por alto el gasto marginal de los reintentos por fallo;
    \item no contar con una estrategia de clasificación de tareas, de modo que incluso las tareas simples pasan por una ruta de búsqueda intensiva.
\end{itemize}
\end{warningbox}

\subsection{Resumen de la sección}
Los resultados experimentales muestran que el desempeño del Agente puede mejorar de forma notable mediante un método sistemático. Pero pasar de la investigación a la puesta en producción todavía requiere optimizar a la vez la tasa de éxito, la latencia y el costo, en lugar de perseguir un solo indicador.
